\subsubsection{Joint consistency check of PP and TG}\label{sec:joint-consistency}
A direct test of shared relaxation fits the PP and all four TG traces simultaneously:
\begin{equation}\label{eq:joint-model}
S_d(t)=B_d+\sum_{j=1}^{m}A_{dj}
\exp\!\left[-\left(\frac{1}{\tau_j}+Dq_d^2\right)t\right],
\qquad q_{\rm PP}=0 .
\end{equation}
The ordered lifetimes $\tau_j$ and nonnegative diffusivity $D$ are shared; each dataset $d$ has its own constant $B_d$ and nonnegative amplitudes $A_{dj}$. This allows reflectivity and diffraction to weight the same relaxation modes differently. The test compares $m=2$ and $m=3$, retaining the original scan exclusions and the \SI{0.5}{ps} starting delay. Initially each trace uses its complete available post-cutoff record. It assumes phase-extracted TG amplitudes, linear optical response, and one common diffusivity for all components.

To prevent the larger PP signal from dominating solely because of its amplitude, define $R_d$ as the measured signal range in the fitted window and minimize
\begin{equation}
\mathcal L=\sum_d\frac{1}{n_d}\sum_{i=1}^{n_d}
\left[\frac{\overline S_d(t_i)-S_d(t_i)}{R_d}\right]^2 .
\end{equation}
Each dataset contributes its squared normalized residual RMS. This is a chosen comparison metric, not a noise-calibrated chi-squared statistic. Linear amplitudes and offsets are solved at each nonlinear step, with multiple starting values for the lifetimes and diffusion coefficient. As a more flexible reference, three exponentials are also fitted independently to each trace. This reference has independent rates, so it need not obey Eq.~\eqref{eq:joint-model}.

The normalized residual RMS values are:
\begin{center}
\begin{tabular}{lrrr}
\toprule
Trace & Joint two-exp. & Joint three-exp. & Independent three-exp.\\
\midrule
PP & 1.639\% & 0.600\% & 0.303\%\\
TG 5.6 & 0.822\% & 0.826\% & 0.681\%\\
TG 6.5 & 0.941\% & 0.728\% & 0.672\%\\
TG 8 & 0.982\% & 0.763\% & 0.692\%\\
TG 12 & 1.323\% & 0.845\% & 0.646\%\\
\bottomrule
\end{tabular}
\end{center}
These RMS values are descriptive diagnostics with unestimated uncertainties. The shared third component improves the total objective and the PP tail, although it does not improve every channel separately. Figures~\ref{fig:joint-pp} and~\ref{fig:joint-tg} show that structured residuals remain, particularly in PP. Independent fits are more flexible by construction; their lower RMS is not by itself a statistical rejection of the shared model. Some independent TG third lifetimes are poorly constrained and approach the broad numerical search limit, so they are not assigned physical meanings.

For uncertainties, 200 bootstrap replicates resample whole retained scans independently within each dataset and refit the shared three-component model. The scan masks and primary normalization factors remain fixed; all replicates converged. To test whether PP's much longer record controls the inference, the entire calculation is repeated using only the common interval from \SI{0.5}{ps} to \SI{26.6}{ps}, with a further 200 successful bootstrap replicates.

\begin{table}[!htbp]
\centering
\caption{Shared three-exponential fit parameters, followed by conditional 95\% whole-scan bootstrap intervals in brackets. The full-record fit uses each dataset's own final delay. The common-window fit truncates every dataset at \SI{26.6}{ps}. Calibration, model, and scan-selection uncertainties are excluded. These parameters describe the assumed joint model, rather than established microscopic lifetimes or material diffusivity.}\label{tab:joint-parameters}
\begin{tabular}{lrr}
\toprule
Parameter & Full records & Common window\\
\midrule
$\tau_1$ (ps) & 0.784 [0.723, 0.844] & 0.528 [0.477, 0.583]\\
$\tau_2$ (ps) & 4.79 [4.30, 5.27] & 1.72 [1.49, 2.18]\\
$\tau_3$ (ps) & 388 [326, 464] & 9.71 [8.56, 13.11]\\
$D$ ($10^{-3}$ m$^2$/s) & 3.80 [2.94, 4.56] & 0.53 [$0^{*}$, 1.81]\\
\bottomrule
\end{tabular}

\smallskip
\parbox{.95\linewidth}{\small $^{*}$The lower percentile reaches the numerical lower bound $D=10^{-8}$ m$^2$/s, shown as zero at the table's precision. The common-window fit therefore does not establish a positive diffusion coefficient away from the boundary.}
\end{table}

The full-record shared model can describe much of both PP and TG with channel-dependent amplitudes. Thus comparing individual biexponential lifetimes was not sufficient to rule out all shared multi-component descriptions: a biexponential approximation can combine the underlying components differently in different channels. Nevertheless, the large changes in Table~\ref{tab:joint-parameters} demonstrate that the shared parameters are not stable against the observation window. Restricting the window removes information about the long PP tail and changes the effective decomposition; the components cannot then be assumed to retain the same physical identities. Removing the $1/n_d$ factor from the objective, while retaining range normalization, changes the full-record solution only slightly; this weighting check does not remove the window dependence.

The stronger test therefore improves empirical consistency without establishing a unique common population model. A third exponential alone does not validate recombination or diffusion assignments. The original slow-TG slope remains the report's conditional baseline estimate, and the joint fit is a separate model check. Raw phase-channel metadata, a controlled matching of acquisition windows, and a justified optical-response model would be needed for a stronger physical interpretation. The calculation and bootstrap samples are saved by \texttt{analysis/joint\_pp\_tg.py} in \texttt{analysis/joint\_pp\_tg.json}.
